%% =============================================================
%%  Curriculum Vitae -- Johora Akter Polin
%%  Compile with pdfLaTeX (Overleaf default); XeLaTeX/LuaLaTeX also work.
%% =============================================================
\documentclass[11pt,letterpaper]{article}

% ---------- Fonts & encoding ----------
\usepackage{iftex}
\usepackage[T1]{fontenc}
\usepackage{tgpagella}            % Palatino-style serif with true small caps
\ifPDFTeX
  \input{glyphtounicode}          % keeps PDF text searchable / ATS-friendly
  \pdfgentounicode=1
\fi
\usepackage{microtype}
\usepackage{fontawesome5}          % LinkedIn logo
\ifPDFTeX
  \newcommand{\scholaricon}{\faGraduationCap}   % pdfLaTeX fallback icon
\else
  \usepackage{academicons}         % Google Scholar logo (XeLaTeX/LuaLaTeX)
  \newcommand{\scholaricon}{\aiGoogleScholar}
\fi

% ---------- Page layout ----------
% Section labels sit in a left column; content sits to the right.
\newlength{\cvlabelwidth}\setlength{\cvlabelwidth}{1.15in}
\newlength{\cvlabelsep}\setlength{\cvlabelsep}{0.2in}
\usepackage[letterpaper,left=2.05in,right=0.7in,top=0.6in,bottom=0.7in,footskip=0.35in]{geometry}

\usepackage{titlesec}
\usepackage{enumitem}
\usepackage{fancyhdr}
\usepackage[hidelinks]{hyperref}  % clickable links, black text, no boxes
\hypersetup{pdftitle={Johora Akter Polin -- Curriculum Vitae},
            pdfauthor={Johora Akter Polin}}

\setlength{\parindent}{0pt}
\setlength{\parskip}{0pt}
\setlength{\emergencystretch}{2em}  % avoids lines running into the margin

% Footer: page number centred on the full page width
\pagestyle{fancy}
\fancyhf{}
\renewcommand{\headrulewidth}{0pt}
\fancyfootoffset[L]{\dimexpr\cvlabelwidth+\cvlabelsep\relax}
\fancyfoot[C]{\small\thepage}

% ---------- Section labels in the left column ----------
\setcounter{secnumdepth}{0}
\titleformat{\section}[leftmargin]
  {\normalfont\large\scshape\filright}{}{0pt}{}
\titlespacing*{\section}{\cvlabelwidth}{2.3ex plus 0.5ex minus 0.3ex}{\cvlabelsep}
% allow a section to start near the page bottom (default keeps the last 20% free)
\renewcommand{\bottomtitlespace}{3\baselineskip}

% ---------- Lists ----------
\setlist{nosep}
\setlist[itemize]{label=\textbullet, leftmargin=1.1em, topsep=2pt, itemsep=1.5pt}
\setlist[enumerate]{labelwidth=2.2em, labelsep=0.5em, leftmargin=2.7em,
                    align=left, topsep=2pt, itemsep=3.5pt}

% ---------- Helper commands ----------
\newcommand{\me}{\textbf{Johora Akter Polin}}
\newcommand{\paperlink}[1]{~\mbox{[\href{#1}{\textsc{paper}}]}}
% bullet line that stays aligned when it is the first thing in a section
\newcommand{\bulletrow}[1]{%
  \hangindent=1.1em\hangafter=1\makebox[1.1em][l]{\textbullet}#1\par\vspace{1.5pt}}
% \cventry{bold heading}{right text}{italic sub-heading}{right text}
\newcommand{\cventry}[4]{%
  \textbf{#1}\hfill #2\par\nobreak
  \textit{#3}\hfill #4\par}
\newcommand{\entrygap}{\par\vspace{5.5pt}}
\newcommand{\pubgroup}[1]{\par\textbf{#1}\par\nobreak}
\newcommand{\pubgap}{\par\vspace{6pt}}
% publication entry: bold title, small citation line, description
\newcommand{\paperitem}[4]{%
  \item[{[#1]}] \textbf{#2.} {\small #3}\par
  \vspace{1pt}#4}
% list style for numbered publication entries ([U1], [C1], ...)
\SetEnumitemKey{publist}{itemsep=3pt, topsep=2pt, leftmargin=2.4em,
  labelwidth=2em, labelsep=0.4em, align=left}
\newcommand{\labeledrow}[2]{%
  \hangindent=1.2em\hangafter=1 \textbf{#1:} #2\par\vspace{2pt}}

% ---------- Header ----------
\newcommand{\cvheader}{%
  \noindent\hspace*{-\dimexpr\cvlabelwidth+\cvlabelsep\relax}%
  \begin{minipage}[b]{\dimexpr\textwidth+\cvlabelwidth+\cvlabelsep\relax}
    \begin{tabular*}{\linewidth}{@{}l@{\extracolsep{\fill}}r@{}}
      {\fontsize{25}{30}\selectfont\scshape Johora Akter Polin} &
      {\small\begin{tabular}[b]{@{}r@{}}
        \href{tel:+17573711873}{\faPhone\ +1 (757) 371-1873}%
        \enspace\textbar\textbar\enspace
        \href{mailto:japolin@wm.edu}{\faEnvelope\ japolin@wm.edu}\\
        \href{https://roboworld314.github.io/japolin.org/}{\faGlobe\ Homepage}%
        \enspace\textbar\textbar\enspace
        \href{https://scholar.google.com/citations?user=50AFRskAAAAJ}{\scholaricon\ Google Scholar}%
        \enspace\textbar\textbar\enspace
        \href{https://www.linkedin.com/in/johora-akter-polin-577759232/}{\faLinkedin\ LinkedIn}
      \end{tabular}}
    \end{tabular*}\par
    \vspace{5pt}\hrule height 0.5pt
  \end{minipage}\par
  \vspace{2pt}}

%% =============================================================
\begin{document}

\cvheader

% -------------------------------------------------------------
\section{Research Interests}
My research builds reliable, efficient, and safety-aware intelligence for
autonomous vehicles and agentic robots. I design perception and
decision-making pipelines that turn camera and LiDAR streams into
timely, trustworthy actions under occlusion, adverse weather, vibration, and
tight compute budgets, and that learn from their failures.
\begin{itemize}[topsep=4pt]
  \item \textbf{Robust Multi-Sensor Perception:} camera-LiDAR fusion,
        occlusion- and blind-zone-aware detection, and all-weather object
        detection for autonomous driving.
  \item \textbf{Agentic AI and Foundation Models:} LLM/VLM reasoning agents
        and multi-agent systems (Qwen, Llama, GPT-4o) combined with
        reinforcement learning, so robots decide at runtime what to perceive,
        retain, and learn from.
  \item \textbf{Edge AI and Real-Time Perception:} latency- and resource-aware
        sensor processing on embedded platforms that keeps decisions anchored
        to fresh sensor evidence.
  \item \textbf{Cooperative and Safety-Critical Autonomy:} V2I-assisted
        collision mitigation, HD-map-based planning, and event-aware Black Box
        recording for failure diagnosis, replay, and data-centric continual learning
        of perception models.
\end{itemize}

% -------------------------------------------------------------
\section{Research Experience}
\cventry{Graduate Research Assistant}{Fall 2025 -- Present}
        {William \& Mary, Williamsburg, VA, USA}{}
Research on perception, agentic AI, and safety for autonomous vehicles with
academic and industry partners, validated in CARLA with ROS~2 and on physical
robotic vehicles (Jetson Orin NX testbed, roadside unit, IEEE 802.11p V2I).
\par\vspace{5pt}
\pubgroup{Publications Under Review}
\begin{itemize}[publist]
  \paperitem{U1}{BlindZoneEye: See Clearly, React Safely, and Record Reliably
    in Autonomous Driving}
    {\me, Yichen Luo, Yongtao Yao, Adwait Nadkarni, and Sidi Lu.
    Under review at an IEEE journal.}
    {Vehicles must act before occluded road users become visible. Fused LiDAR
    geometry, raycasting visibility reasoning, and YOLOv11 camera semantics
    into a blind-zone risk score driving V2I- and HD-map-assisted planning and an event-aware Black Box (92.3\% detection
    accuracy at 41\,ms; 76.2\% fewer collisions).}
  \paperitem{U2}{SenseGuard: Sensor-to-Decision Adaptation and Learning for
    Reliable Agentic Robots}
    {\me, Yichen Luo, Yongtao Yao, Adwait Nadkarni, and Sidi Lu.
    Under review at an international robotics conference.}
    {Agentic robots act on stale inputs when redundant frames clog
    pipelines. Combined SigLIP similarity, Qwen2.5-VL chain-of-thought
    reasoning, and a PPO policy to filter camera and LiDAR frames and keep
    failure episodes for retraining (up to +20 mAP@50, 31\% lower latency,
    75\% lower LiDAR trajectory error).}
  \paperitem{U3}{VibraBox: A Vibration-Adaptive Black Box for Replay and
    Learning from Failures in Mobile-Robot Autonomy}
    {\me, Yichen Luo, Yongtao Yao, and Sidi Lu.
    Under review at an international robotics conference.}
    {Robot failures are hard to replay, and terrain vibration changes how much
    LiDAR data can safely be dropped. Built a LiDAR Black Box with
    BEV-histogram (Chi-square) deduplication, voxel-grid downsampling, and
    LZ4/ZSTD compression, tuned by a DQN controller from vibration and KISS-ICP
    stability (up to 46\% less storage, 71.4\% lower rotational replay error,
    +12.7\% recall after retraining).}
\end{itemize}
\pubgap
\pubgroup{Conference Papers}
\begin{itemize}[publist]
  \paperitem{C1}{Log What Matters: Safety-Aware Adaptive Logging for CAVs in
    Dynamic Conditions}
    {\me, Yichen Luo, FNU Sumaiya, Zheng Dong, and Sidi Lu. \textit{2026 IEEE
    4th International Conference on Mobility, Operations, Services and
    Technologies (MOST)}, Detroit, MI, USA, May
    2026.\paperlink{https://sidilu.org/paper/AdaptiveLogAsInput_MOST26.pdf}}
    {Redundant frames waste compute and force vehicles to act on stale
    perception. Proposed CAAL-VLM, adaptive camera and LiDAR logging
    using SigLIP, Qwen-VL, and RL (PPO, DQN) policies, plus GPT-4o sensor
    selection (27\% less LiDAR storage, 42\% lower latency, rainy mAP from
    0.75 to 0.86).}
\end{itemize}
\pubgap
\pubgroup{Workshop Papers}
\begin{itemize}[publist]
  \paperitem{W1}{Context-Aware Perception: VLM-Augmented All-Weather Detection
    for Autonomous Driving}
    {\me, Yichen Luo, and Sidi Lu. \textit{The Tenth ACM/IEEE Symposium on
    Edge Computing (SEC'25)}, Arlington, VA, USA,
    2025.\paperlink{https://doi.org/10.1145/3769102.3774638}}
    {Camera detectors lose accuracy in rain, fog, and snow. Used
    vision-language models to extract weather and visibility context from
    driving scenes and adapt object detection for robust all-weather
    autonomous driving.}
  \paperitem{W2}{TinyLite: Efficient Edge Intelligence for Real-Time Spam
    Detection Under Device Constraints}
    {Xuanlin Zhu, \me, and Sidi Lu. \textit{IEEE Workshop on Physical
    Intelligence: Systems and Applications (PISA), at IEEE MOST}, May
    2026.\paperlink{https://sidilu.org/paper/PISA_26_SpamDetection.pdf}}
    {Contributed to TinyLite, a compressed BERT-Tiny Transformer for on-device
    spam detection that cuts FLOPs by 81.6\% with under 1.1 points of F1 loss.}
\end{itemize}
\entrygap
\cventry{Machine Learning Research}{2020 -- 2024}
        {Daffodil International University, Dhaka, Bangladesh}{}
Applied CNNs and supervised machine learning to agriculture, natural-language
processing, social-impact, and security problems.
\par\vspace{5pt}
\pubgroup{Conference Papers}
\begin{itemize}[publist]
  \paperitem{C2}{Encryption Process of Blockchain-Based Online Course
    Curriculum Education System}
    {Md Rayhan Ahmmed, Md Mamun Sakib, Maruf Ur Rahman, Rabeya Bibi,
    Md Maruf Hasan Talukder, and \me. \textit{11th International Symposium on
    Digital Forensics and Security (ISDFS)},
    2023.\paperlink{https://ieeexplore.ieee.org/abstract/document/10131770}}
    {}
  \paperitem{C3}{Predicting Sugarcane Yields using Supervised Learning: A
    Comparative Study}
    {Parvez Ahmed Papon, Md Mahfuzur Rahman, Shafin Ahmed, Shahriar Mamun,
    Md Zahirul Islam Mehadi, and \me. \textit{14th International Conference
    on Computing Communication and Networking Technologies (ICCCNT)},
    2023.\paperlink{https://ieeexplore.ieee.org/abstract/document/10307590}}
    {}
  \paperitem{C4}{Utilizing Smartphone Usage Pattern for Gender Recognition
    Applying Machine Learning Algorithms}
    {\me, Omayer Khan, and Ahmed Al Marouf. \textit{4th International
    Conference on Electronics, Communication and Aerospace Technology
    (ICECA)}, 2020.\paperlink{https://ieeexplore.ieee.org/abstract/document/9297407}}
    {}
\end{itemize}
\pubgap
\pubgroup{Journal Articles}
\begin{itemize}[publist]
  \paperitem{J1}{A Novel Approach to Analyzing the Impact of AI, ChatGPT, and
    Chatbot on Education Using Machine Learning Algorithms}
    {Nahid Hasan, \me, Md.~Rayhan Ahmmed, Md.~Mamun Sakib, Md.~Farhan Jahin,
    and Md.~Mahfuzur Rahman.
    \textit{Bulletin of Electrical Engineering and Informatics}, vol.~13, no.~4, pp.~2951--2958,
    2024.\paperlink{https://doi.org/10.11591/eei.v13i4.7158}}
    {}
  \paperitem{J2}{Predicting the Effects of Microcredit on Women's Empowerment
    in Rural Bangladesh: Using Machine Learning Algorithms}
    {\me, Md.~Fouad Hossain Sarker, MST Dilruba Khanom Dolon, Nahid Hasan,
    Md.~Mahfuzur Rahman, and Zannatun Nayem Vasha. \textit{Bulletin of
    Electrical Engineering and Informatics}, vol.~13, no.~4, pp.~2699--2706,
    2024.\paperlink{https://doi.org/10.11591/eei.v13i4.7169}}
    {}
  \paperitem{J3}{Tomato Pest Recognition Using Convolutional Neural Network in
    Bangladesh}
    {\me, Nahid Hasan, Md.~Tarek Habib, Atiqur Rahman, Zannatun Nayem Vasha,
    and Bidyut Sharma. \textit{Bulletin of Electrical Engineering and
    Informatics}, vol.~13, no.~1, pp.~619--627, 2024.\paperlink{https://doi.org/10.11591/eei.v13i1.6073}}
    {}
  \paperitem{J4}{Bengali Slang Detection Using State-of-the-Art Supervised
    Models from a Given Text}
    {Md Abdul Hamid, Eteka Sultana Tumpa, \me, Jabir Al Nahian, Atiqur
    Rahman, and Nurjahan Akther Mim. \textit{Bulletin of Electrical
    Engineering and Informatics}, vol.~12, no.~4, pp.~2381--2387,
    2023.\paperlink{https://doi.org/10.11591/eei.v12i4.4743}}
    {}
  \paperitem{J5}{Depression Detection in Social Media Comments Data Using
    Machine Learning Algorithms}
    {Zannatun Nayem Vasha, Bidyut Sharma, Israt Jahan Esha, Jabir Al Nahian,
    and \me. \textit{Bulletin of Electrical Engineering and Informatics}, vol.~12, no.~2, pp.~987--996,
    2023.\paperlink{https://doi.org/10.11591/eei.v12i2.4182}}
    {}
\end{itemize}

% -------------------------------------------------------------
\section{Education}
\cventry{William \& Mary}{Williamsburg, VA, USA}
        {Ph.D. in Computer Science}{2024 -- 2029 (expected)}
\begin{itemize}
  \item Advisor: Dr.~Sidi Lu; GPA: 3.95/4.00
  \item Research Area: Edge and Vehicle Computing for Autonomous Driving
\end{itemize}
\entrygap
\cventry{Daffodil International University}{Dhaka, Bangladesh}
        {M.Sc. in Computer Science and Engineering}{2021 -- 2023}
\begin{itemize}
  \item Advisor: Dr.~Md.~Tarek Habib; GPA: 4.00/4.00
\end{itemize}
\entrygap
\cventry{Daffodil International University}{Dhaka, Bangladesh}
        {B.Sc. in Computer Science and Engineering}{2015 -- 2020}
\begin{itemize}
  \item GPA: 3.90/4.00
\end{itemize}
\entrygap
\cventry{HEC Montr\'eal}{Montr\'eal, Canada}
        {Certification Course in Social Business (3 credits)}{2020}

% -------------------------------------------------------------
\section{Technical Skills}
\labeledrow{Agentic AI \& Foundation Models}{Qwen family (Qwen2.5-VL,
  Qwen2.5, Qwen3, Qwen-Agent), Llama, GPT-4o (OpenAI API), SigLIP;
  VLM chain-of-thought reasoning agents, multi-agent systems, and
  reinforcement-learning agents (PPO, DQN); model serving with vLLM and Ollama}
\labeledrow{Programming}{Python, C, C\#, MATLAB; PyTorch}
\labeledrow{Perception \& ML}{Hands-on experience with 25+ camera-based
  perception models, including the YOLO family (v8--v11), RF-DETR, RT-DETR,
  DINO, and Faster/Cascade R-CNN; camera-LiDAR fusion; CNNs}
\labeledrow{LiDAR \& 3D Perception}{64- and 128-beam 360\textdegree{} LiDAR
  (vehicle and roadside); point-cloud processing (BEV projection, voxel-grid
  filtering, farthest point sampling, occupancy grids, raycasting); 3D
  detection (CenterPoint, PointPillars, SECOND); LiDAR odometry with KISS-ICP
  and trajectory evaluation (ATE, ARE, RPE) against GNSS/INS ground truth;
  lossless point-cloud compression (LZ4, ZSTD, G-PCC)}
\labeledrow{Autonomy \& Robotics}{ROS~2 (rosbag2/MCAP logging), CARLA, HD maps
  (OpenDRIVE), model predictive control, V2I communication (IEEE 802.11p,
  DDS/RTPS)}
\labeledrow{Platforms}{NVIDIA Jetson Orin NX/Nano; Yahboom ROSMASTER X3 and R2
  robots; Robify RobiX-Nav and RobiS autonomous driving robots; roadside unit
  (RSU)}
\labeledrow{Languages}{English (IELTS overall band score: 7.0), Bangla}

% -------------------------------------------------------------
\section{Teaching \& Mentoring}
\cventry{Assistant Instructor \textnormal{(\href{https://roboworld314.github.io/VC_Lab/}{roboworld314.github.io/VC\_Lab})}}{Fall 2026}
        {William \& Mary, Williamsburg, VA, USA}{}
\begin{itemize}
  \item Assisted Dr.~Sidi Lu in teaching \textit{COLL/CSCI 100: Physical
        Intelligence: Robots}, a hands-on course in which students program
        5G-connected robotic vehicles on a reconfigurable smart-city roadway
        with programmable traffic lights.
  \item Supervised project milestones spanning CARLA smart-city scenario
        design, camera and LiDAR data collection with YOLO object detection,
        safety-aware robot responses to traffic signals and obstacles, and
        digital-twin HD mapping.
  \item Delivered guest lectures on physical intelligence, autonomous vehicles,
        and edge computing, with live demonstrations on robotic vehicle
        platforms.
  \item Created and maintain the course website.
\end{itemize}
\entrygap
\cventry{Graduate Teaching Assistant}{Fall 2024 -- Spring 2025}
        {William \& Mary, Williamsburg, VA, USA}{}
\begin{itemize}
  \item Conducted programming lab sessions and advised students on coursework.
\end{itemize}
\entrygap
\cventry{Senior Lecturer}{July 2023 -- August 2024}
        {Daffodil International University, Dhaka, Bangladesh}{}
\begin{itemize}
  \item Taught theory and laboratory courses in Computer Science and Engineering.
  \item Collaborated with students on research projects.
\end{itemize}
\entrygap
\cventry{Lecturer}{January 2021 -- June 2023}
        {Daffodil International University, Dhaka, Bangladesh}{}
\begin{itemize}
  \item Conducted theory and lab sessions; evaluated, advised, and tracked students.
\end{itemize}
\entrygap
\cventry{Judge \& Mentor}{January 2021 -- September 2021}
        {Social Business Creation Competition, HEC Montr\'eal, Canada}{}
\begin{itemize}
  \item Served on the judging committee evaluating student social-business
        proposals.
  \item Mentored participating teams on business-model development and
        scale-up.
\end{itemize}

% -------------------------------------------------------------
\section{Awards \& Honors}
\bulletrow{Student Travel Grant, IEEE MOST 2026}
\bulletrow{Student Travel Grant, ACM/IEEE Symposium on Edge Computing (SEC) 2025}
\bulletrow{Recognized as Advisor, Girls Computer Programming Club, 2022}
\bulletrow{Core Technical Member, National Hackathon, CTO Forum Bangladesh, 2021}
\bulletrow{Judging Committee Member, Social Business Creation Competition
  (Business Model Innovation and Execution \& Scale-up rounds), 2020--2022}
\bulletrow{Scopus-Indexed Publication Recognition, Daffodil International
  University, 2020}
\bulletrow{Runner-up Team, Social Business Creation Global Final, HEC
  Montr\'eal, 2019}

% -------------------------------------------------------------
\section{Academic Service}
\labeledrow{Session Chair}{1st International Workshop on Edge Computing for
  Connected and Autonomous Vehicles (EdgeCAV'25), at SEC'25}
\labeledrow{Reviewer}{EdgeCAV'25 and EdgeCAV'26 workshops; ACM/IEEE Symposium
  on Edge Computing (SEC'25, SEC'26)}
\labeledrow{Panelist}{Student/Alumni Panel (autonomous driving research),
  Computer Science Symposium for Graduate Studies, William \& Mary,
  November 2026}
\labeledrow{Volunteer}{IEEE MOST'25 Conference; CS Symposium, 2024}
\labeledrow{Organizer}{National Girls Programming Contest; National Hackathon}
\labeledrow{Participant}{\href{https://capwic.org/conferences/2025/}{ACM
  CAPWIC (Celebration of Women in Computing)}, 2025; Workshop
  on Developing Quality Research, 2023; Social Business Program, Canada}

\end{document}
